% ============================================================================
% Fichier  : partie6/P06_C29_protocole_ZKNR.tex
% Titre    : Le Protocole ZK-NR -- Architecture et Statut d'un Travail
%            de Recherche en Cours
% Auteur   : MINKA MI NGUIDJOÏ Thierry Emmanuel
% Version  : 2.0 -- Juin 2026
% ============================================================================

\chapter{Le Protocole ZK-NR. Architecture et Statut d'un Travail
de Recherche en Cours}
\label{chap:ZKNR}

\epigraph{%
    « Zero-knowledge proofs represent one of the most powerful tools in
    cryptography, verification without disclosure, truth without
    exposure. »%
}{Shafi Goldwasser}

\niveauChapitre{Recherche}{%
    Preuves à divulgation nulle de connaissance (Ch.~\ref{chap:CRO},
    \ref{chap:prim-crypt}) recommandées. Chapitre~\ref{chap:cryptanal}
    (cryptanalyse, prudence sur les travaux en cours) indispensable~:
    ce chapitre applique scrupuleusement la même discipline.%
}

\begin{boxobjectifs}
\begin{itemize}
    \item Décrire l'objectif annoncé du protocole ZK-NR
          (\textit{Zero-Knowledge Non-Repudiation}) et son architecture
          de principe, déductible des primitives nommées.
    \item Situer chacune des quatre primitives composantes
          (Merkle, STARK, signatures à seuil, signature post-quantique)
          dans le cadre théorique déjà établi (Ch.~\ref{chap:CRO}
          à~\ref{chap:quantique}).
    \item Comprendre précisément ce que signifie le \textbf{statut
          de prépublication} en recherche scientifique, et pourquoi
          il impose une lecture critique distincte de celle d'un
          article publié et évalué par les pairs.
    \item Appliquer une grille de lecture critique à une annonce de
          protocole de recherche, compétence transférable à
          toute situation où un investigateur rencontre un outil
          ou un mécanisme cryptographique émergent, non encore
          validé par la communauté scientifique.
    \item Distinguer ce qu'une description architecturale permet
          d'affirmer en toute rigueur de ce qui resterait à établir
          par une publication complète et évaluée par les pairs.
\end{itemize}
\end{boxobjectifs}

\bigskip

% ============================================================================
\section*{Un chapitre sur un protocole, et sur la lecture critique
d'un protocole}
% ============================================================================

Ce chapitre porte sur ZK-NR (\textit{Zero-Knowledge Non-Repudiation}),
un protocole de recherche de l'auteur de ce cours, décrit dans une
prépublication déposée sur l'IACR ePrint~\cite{minka2025zknr}, conçu
pour la non-répudiation préservant la vie privée avec des garanties
post-quantiques. Mais il porte, tout autant, sur une question plus
générale et directement utile à votre pratique
professionnelle~: \textbf{comment lire et évaluer rigoureusement
l'annonce d'un protocole cryptographique qui n'existe encore qu'en
prépublication, sans évaluation par les pairs}, une situation que vous
rencontrerez à de nombreuses reprises, face à des outils
commerciaux, des préprints académiques, ou des annonces
technologiques, tout au long de votre carrière d'investigateur.

\separateur

% ============================================================================
\section{Objectif Annoncé et Architecture de Principe}
\label{sec:objectif-architecture}
% ============================================================================

\subsection{L'objectif visé}

ZK-NR vise à répondre à un problème directement issu du Trilemme CRO
formalisé au Ch.~\ref{chap:CRO}~: produire des attestations
numériques offrant \textbf{simultanément}~:
\begin{itemize}
    \item une \textbf{non-répudiation forte}. l'émetteur d'une
          attestation ne peut nier l'avoir produite, propriété
          centrale de l'opposabilité juridique
          (Ch.~8~(UE~SEC4031), Ch.~\ref{chap:rapport})~;
    \item une \textbf{préservation de la confidentialité}.
          ne révéler que le strict nécessaire, sur le modèle des
          preuves à divulgation nulle (Ch.~\ref{chap:CRO},
          Section~4.4)~;
    \item une \textbf{résistance post-quantique}. anticiper,
          conformément au Ch.~\ref{chap:quantique}, la menace que
          représente l'algorithme de Shor pour les mécanismes
          de signature classiques.
\end{itemize}

\begin{boxCRO}
\textbf{Analyse CRO. ZK-NR comme tentative architecturale face
au Trilemme}

Le Théorème~3 (Ch.~\ref{chap:CRO}) établit qu'aucune primitive
unique ne peut maximiser simultanément $\mathsf{Conf}$,
$\mathsf{Rel}$ et $\mathsf{Opp}$. L'objectif affiché de ZK-NR
illustre précisément la \textbf{réponse architecturale} que ce
théorème appelle~: non pas chercher une primitive miracle, mais
\textbf{combiner plusieurs primitives complémentaires}, chacune
apportant sa force dominante, exactement la logique de
l'architecture Gold-Iron-Clay déjà présentée comme optimale dans
sa classe de construction (Ch.~\ref{chap:CRO}, Section~6.1).

\cro{$\sim$ objectif visé}{$\sim$ objectif visé}{$\sim$ objectif visé}

\textbf{Note de prudence}~: cette analyse porte sur la
\textbf{cohérence de l'objectif annoncé} avec le cadre théorique
établi, pas sur une évaluation de l'atteinte effective de cet
objectif par l'implémentation, question qui ne peut être
tranchée qu'à l'examen d'une publication complète.
\end{boxCRO}

\subsection{Architecture de principe~: quatre primitives, quatre rôles}

L'architecture annoncée combine quatre couches, chacune correspondant
à une primitive \textbf{individuellement réelle et bien documentée}
dans la littérature scientifique, ce qui permet d'en discuter
le rôle de principe sans dépendre d'une publication ZK-NR complète.

\begin{tableaucours}{p{2.8cm}p{3.2cm}p{5.0cm}}{%
    \tableauHeader{Couche} &
    \tableauHeader{Primitive sous-jacente} &
    \tableauHeader{Rôle de principe (déductible des propriétés
    publiques de la primitive)}
}
\textbf{Engagement} & Arbre de Merkle &
    Structure d'engagement immuable~: garantit qu'un ensemble de
    données ne peut être modifié sans que la modification soit
    détectable, principe déjà rencontré au Ch.~\ref{chap:for-sys}
    pour l'intégrité de structures de données \\
\textbf{Preuve} & STARK &
    Preuve à divulgation nulle de connaissance, \textbf{transparente}
    (sans \textit{trusted setup}), propriété qui la distingue
    structurellement des zk-SNARKs analysés au
    Ch.~\ref{chap:prim-crypt}, lesquels nécessitent une configuration
    de confiance initiale \\
\textbf{Signature} & Signature à seuil (BLS ou
    équivalent) & Permet une attestation collective nécessitant
    l'accord d'un sous-ensemble de signataires, analysée en
    détail au Ch.~\ref{chap:cas-zknr} (vulnérabilité face à Shor,
    Tableau~32.1) \\
\textbf{Authentification} & Signature post-quantique
    (Dilithium/ML-DSA ou équivalent) & Résistance à long terme face
    à la menace quantique (Ch.~\ref{chap:quantique}, FIPS~204),
    compense, en principe, la vulnérabilité de la couche précédente \\
\end{tableaucours}
\vspace{-0.8em}
\begin{center}{\small\itshape Tableau~29.1. Architecture de principe de
ZK-NR, primitive par primitive}\end{center}

\begin{boiteRemarque}{Ce que ce tableau établit, et ce qu'il
n'établit pas}
Chaque ligne du Tableau~29.1 décrit une propriété \textbf{de la
primitive elle-même}, vérifiable indépendamment dans la littérature
citée tout au long de ce cours (Merkle~: Ch.~\ref{chap:for-sys}~;
STARK~: Ch.~\ref{chap:quantique}, Ch.~\ref{chap:prim-crypt}~;
BLS~: Ch.~\ref{chap:cas-zknr}~; Dilithium~: Ch.~\ref{chap:quantique},
FIPS~204). Ce tableau \textbf{n'établit pas}, et ne peut pas
établir, en l'absence d'article complet, que leur
\textbf{combinaison spécifique} dans ZK-NR atteint effectivement
les propriétés de sécurité visées, ni quelle est la complexité
exacte de leur orchestration, ni quelles garanties formelles ont
été démontrées pour l'ensemble du système plutôt que pour chaque
composant pris isolément.
\end{boiteRemarque}

% ============================================================================
\section{Sur le Statut de Prépublication en Recherche Scientifique}
\label{sec:statut-prepublication}
% ============================================================================

\subsection{Une distinction essentielle}

\begin{boxdef}
\textbf{Prépublication}\index{prépublication}
(\textit{preprint})~: document de recherche rendu public par ses
auteurs, typiquement sur une plateforme telle qu'\textit{ePrint}
(IACR) ou \textit{arXiv}, \textbf{avant} ou \textbf{indépendamment}
d'une évaluation formelle par les pairs (\textit{peer review}) dans
le cadre d'une revue ou conférence scientifique.

\textbf{Publication évaluée par les pairs}~: document ayant subi
un processus d'évaluation indépendant par des experts du domaine,
mandatés par une revue ou conférence, vérifiant la rigueur des
démonstrations, la nouveauté de la contribution, et la validité
des conclusions, avant acceptation et publication officielle.
\end{boxdef}

\begin{boxalert}
\textbf{La distinction n'est pas une question de qualité présumée,
mais de statut de vérification}

Une prépublication n'est \textbf{pas nécessairement de moindre
qualité} qu'une publication évaluée par les pairs, de nombreux
résultats majeurs de l'histoire des mathématiques et de
l'informatique ont d'abord circulé en prépublication. La
différence est d'une autre nature~: une prépublication
\textbf{n'a pas encore été soumise au contrôle indépendant} qui
permet à un tiers de s'appuyer dessus avec la même confiance
qu'un résultat publié et vérifié. C'est exactement la distinction
qui sépare, dans ce cours, le statut du Trilemme CRO \textbf{avant}
sa publication dans IEEE Access (Ch.~\ref{chap:CRO}, où il aurait
fallu le traiter avec cette même prudence) de son statut
\textbf{après} cette publication, où les Théorèmes~1 à~3 peuvent
être cités et utilisés comme des résultats établis.
\end{boxalert}

\subsection{Application au cas ZK-NR}

\begin{tableaucours}{p{3.5cm}p{4.0cm}p{4.5cm}}{%
    \tableauHeader{Critère} &
    \tableauHeader{Trilemme CRO (Ch.~\ref{chap:CRO}-\ref{chap:prim-crypt})} &
    \tableauHeader{ZK-NR (ce chapitre)}
}
Statut de publication & Publié et évalué par les pairs (IEEE Access,
    2026) & Prépublication sur l'IACR ePrint
    (2025/1138)~\cite{minka2025zknr}~; non encore évaluée par les
    pairs ni auditée \\
Démonstrations & Théorèmes~1, 2, 3 entièrement démontrés
    et vérifiés (Pinsker, Fano) & Architecture de principe
    décrite (Tableau~29.1)~; preuves de sécurité formelle non
    accessibles pour vérification \\
Validation empirique & Trois études de cas calculées et publiées
    (Helios, zk-SNARK, ECDSA) & Aucune donnée chiffrée publiée
    et vérifiable \\
\textbf{Usage approprié dans ce cours} & Citer les théorèmes,
    reproduire les calculs, s'appuyer sur les résultats
    & Décrire l'objectif et l'architecture de principe~;
    \textbf{ne pas} affirmer de propriétés de sécurité
    spécifiques non vérifiées \\
\end{tableaucours}
\vspace{-0.8em}
\begin{center}{\small\itshape Tableau~29.2. Comparaison de statut~:
Trilemme CRO publié vs.\ ZK-NR en développement}\end{center}

% ============================================================================
\section{Grille de Lecture Critique d'une Annonce de Protocole}
\label{sec:grille-lecture}
% ============================================================================

Cette section généralise la démarche appliquée à ZK-NR en une
grille transférable, directement utile face à \textbf{tout}
protocole ou outil cryptographique émergent rencontré dans votre
pratique professionnelle, qu'il s'agisse d'une solution
commerciale présentée par un fournisseur, d'un préprint académique,
ou d'une annonce technologique relayée par la presse spécialisée.

\begin{tableaucours}{p{3.2cm}p{8.8cm}}{%
    \tableauHeader{Question} &
    \tableauHeader{Ce qu'elle permet de vérifier}
}
\textbf{Quel est le statut de
    publication exact~?} & Prépublication (ePrint, arXiv)~?
    Publication évaluée par les pairs (revue, conférence)~? Simple
    communiqué de presse ou documentation commerciale sans
    publication scientifique sous-jacente~? \\
\textbf{Les primitives composantes
    sont-elles individuellement
    documentées~?} & Chaque brique cryptographique utilisée
    a-t-elle, elle-même, fait l'objet d'une analyse de sécurité
    publiée et reconnue (sur le modèle du Tableau~29.1 pour ZK-NR)~? \\
\textbf{Les propriétés de sécurité
    globales sont-elles
    démontrées, ou seulement
    affirmées~?} & Existe-t-il une preuve formelle complète
    (Ch.~\ref{chap:form-meth}) de la combinaison, ou seulement une
    description qualitative des propriétés espérées~? \\
\textbf{Des scores numériques
    précis sont-ils avancés~?} & Si oui, une méthode de calcul
    complète est-elle publiée et reproductible (sur le modèle du
    Ch.~\ref{chap:prim-crypt})~? Un score sans méthode vérifiable
    doit être traité avec le plus grand scepticisme. \\
\textbf{Quelles limitations les
    auteurs reconnaissent-ils
    eux-mêmes~?} & Les auteurs sérieux identifient explicitement
    les limites de leur propre travail (cf.~le Trilemme CRO,
    Ch.~\ref{chap:CRO}, qui qualifie ses propres bornes de
    ``open problem'' sur certains aspects)~; l'absence totale de
    limitation reconnue est elle-même un signal d'alerte. \\
\end{tableaucours}
\vspace{-0.8em}
\begin{center}{\small\itshape Tableau~29.3. Grille de lecture critique
d'une annonce de protocole cryptographique}\end{center}

\begin{boxjuris}
\textbf{Application judiciaire de cette grille~: le témoignage
d'expert face à un protocole émergent}

Un expert judiciaire camerounais (Ch.~12~(UE~SEC4031)) appelé
à se prononcer sur la fiabilité d'un mécanisme cryptographique
récent, utilisé, par exemple, par une plateforme Mobile Money
pour sécuriser ses transactions (Ch.~\ref{chap:for-mob}), doit
appliquer cette même grille avant de fonder une conclusion sur ce
mécanisme. Affirmer la fiabilité d'un protocole sur la seule base
de la documentation commerciale du fournisseur, sans vérifier son
statut de publication et de validation indépendante
(Tableau~29.3), exposerait l'expert à une contestation directe en
contre-interrogatoire (Ch.~\ref{chap:rapport}), exactement le
type de question illustrée à l'Exercice~32.3 de ce cours
concernant le protocole fictif ``CamSeal''.
\end{boxjuris}

% ============================================================================
\section{Ce que ZK-NR Illustre pour la Pratique Forensique}
\label{sec:illustration-pratique}
% ============================================================================

Indépendamment du statut de publication de ZK-NR lui-même, son
\textbf{architecture de principe} illustre des leçons directement
applicables à la pratique forensique de ce cours.

\begin{enumerate}
    \item \textbf{La non-répudiation post-quantique est un problème
          réel et actuel}~: le Ch.~\ref{chap:quantique} a établi
          que les signatures de la chaîne de custody numérique
          d'aujourd'hui (Ch.~8~(UE~SEC4031)) pourraient perdre
          leur valeur probatoire si l'algorithme sous-jacent
          devient cassable. Tout effort de recherche visant une
          non-répudiation résistante au quantique répond à un
          besoin identifié, indépendamment de l'aboutissement
          spécifique de tel ou tel protocole.
    \item \textbf{La préservation de la vie privée et l'opposabilité
          judiciaire sont en tension structurelle}~: le Théorème~3
          (Ch.~\ref{chap:CRO}) l'a démontré mathématiquement. Tout
          protocole visant à combiner les deux doit nécessairement
          composer plusieurs primitives, aucune primitive unique
          ne peut s'y soustraire.
    \item \textbf{La recherche en cryptographie appliquée à la
          forensique est un domaine actif}~: ce cours lui-même,
          par ses six chapitres de Partie~VI, illustre que les
          questions soulevées par la pratique forensique
          camerounaise (Ch.~12~(UE~SEC4031),
          Ch.~\ref{chap:bench}) trouvent un écho direct dans la
          recherche scientifique internationale en cryptographie
          institutionnelle.
\end{enumerate}

\separateur

\begin{boxresume}
\textbf{Ce qu'il faut retenir de ce chapitre~:}
\begin{itemize}
    \item \textbf{ZK-NR} vise une non-répudiation préservant la vie
          privée avec garanties post-quantiques, en combinant quatre
          primitives~: engagement de Merkle, preuve STARK
          (transparente), signature à seuil, signature
          post-quantique, chaque primitive étant individuellement
          réelle et documentée, indépendamment du statut du protocole
          combiné.

    \item Cette architecture illustre, par sa conception même, la
          \textbf{réponse au Trilemme CRO} (Ch.~\ref{chap:CRO})~:
          combiner plusieurs primitives complémentaires plutôt que
          chercher une primitive unique impossible.

    \item \textbf{Distinction essentielle}~: \textbf{prépublication}
          (statut non vérifié par les pairs) vs.\ \textbf{publication
          évaluée par les pairs} (comme le Trilemme CRO lui-même,
          Ch.~\ref{chap:CRO}-\ref{chap:prim-crypt}). Cette distinction
          détermine ce qui peut être affirmé avec assurance et ce
          qui doit être traité avec prudence.

    \item \textbf{Grille de lecture critique en cinq questions}
          (Tableau~29.3), transférable à \textbf{tout} protocole
          ou outil cryptographique émergent rencontré dans la
          pratique professionnelle~: statut de publication,
          documentation des primitives composantes, preuves
          formelles vs.\ affirmations qualitatives, vérifiabilité
          des scores numériques, limitations reconnues par les
          auteurs.

    \item \textbf{Application judiciaire directe}~: un expert
          appelé à témoigner sur la fiabilité d'un mécanisme
          cryptographique émergent doit appliquer cette grille
          avant toute conclusion, sous peine d'exposer son
          témoignage à une contestation fondée en
          contre-interrogatoire.

    \item Ce chapitre lui-même constitue une \textbf{application de
          sa propre leçon}~: il décrit ce qui peut être établi sur
          ZK-NR sans franchir la limite vers des affirmations non
          vérifiables, exactement la discipline exigée de tout
          rapport forensique (Ch.~\ref{chap:rapport}).
\end{itemize}
\end{boxresume}

\bigskip

\begin{boxexo}[title={\ding{45}\quad Exercice~29.1. Distinguer les statuts \niveaubase}]
Pour chacune des affirmations suivantes, indiquez si elle relève
d'un résultat \textbf{établi} (publié et vérifié), d'une
\textbf{architecture de principe} (déductible des propriétés
publiques des composants), ou d'une \textbf{affirmation non
vérifiable} en l'état des sources de ce cours~:
\begin{enumerate}[(a)]
    \item ``Les STARKs ne nécessitent pas de \textit{trusted
          setup}.''
    \item ``ZK-NR atteint un score CRO de~$0{,}88$ pour
          l'opposabilité.''
    \item ``BLS est vulnérable à l'algorithme de Shor.''
    \item ``Le modèle Tamarin de ZK-NR comporte huit règles
          symboliques et prouve six lemmes critiques.''
    \item ``Le Trilemme CRO conditionnel établit
          $\mathsf{Conf}\cdot\mathsf{Rel}\cdot\mathsf{Opp} \leq
          \Delta/\lambda + O(\lambda^{-2})$.''
\end{enumerate}
\end{boxexo}

\begin{boxexo}[title={\ding{45}\quad Exercice~29.2. Appliquer la grille de lecture \niveaumoyen}]
Une entreprise présente sur son site web un nouveau service de
``signature électronique quantique-sûre'' pour les contrats
commerciaux, sans référence à une publication scientifique, mais
avec un argumentaire technique détaillé citant ``une architecture
propriétaire combinant plusieurs couches de sécurité''.
\begin{enumerate}[(a)]
    \item Appliquez les cinq questions du Tableau~29.3 à cette
          situation. Pour chacune, formulez la question précise
          que vous poseriez à cette entreprise.
    \item En quoi cette situation diffère-t-elle, sur le plan du
          statut de vérifiabilité, du cas ZK-NR traité dans ce
          chapitre~? En quoi est-elle similaire~?
    \item Si cette entreprise sollicitait une attestation de votre
          part, en tant qu'expert, sur la fiabilité de son
          mécanisme pour un usage dans des contrats à enjeu élevé,
          quelle serait votre position, conformément à la
          discipline de ce chapitre~?
\end{enumerate}
\end{boxexo}

\begin{boxexo}[title={\ding{45}\quad Exercice~29.3. Rédiger une fiche de veille technologique \niveauexpert}]
Vous êtes chargé, au sein d'une cellule de veille technologique de
l'ANTIC, de produire une fiche de synthèse sur le protocole ZK-NR
pour informer les décideurs institutionnels camerounais, sans
exagérer ni minimiser son potentiel.
\begin{enumerate}[(a)]
    \item Rédigez une section ``Objectif et architecture'' (5--8
          lignes), en vous appuyant exclusivement sur ce qui est
          établi dans ce chapitre (Tableau~29.1), sans ajouter
          d'affirmations non sourcées.
    \item Rédigez une section ``Statut et limitations'' (5--8
          lignes), appliquant la distinction prépublication~/
          publication évaluée par les pairs
          (Section~\ref{sec:statut-prepublication}).
    \item Rédigez une section ``Pertinence pour le contexte
          camerounais'' (5--8 lignes), reliant l'objectif de
          ZK-NR aux besoins identifiés dans ce cours (chaîne de
          custody post-quantique, Ch.~\ref{chap:quantique}~;
          formation des magistrats face à des preuves
          cryptographiques complexes, Ch.~\ref{chap:bench}).
    \item Concluez par une recommandation~: à quel horizon
          temporel et sous quelles conditions ce type de protocole
          pourrait-il, selon vous, être considéré pour une
          adoption institutionnelle au Cameroun~? Justifiez en
          vous appuyant sur les régimes de capacité institutionnelle
          du Ch.~\ref{chap:prim-crypt}.
\end{enumerate}
\end{boxexo}

\bigskip

\begin{boxinfo}
\textbf{Transition.} Ce chapitre a complété la présentation du
cadre théorique de la Partie~VI en abordant directement le
protocole ZK-NR, avec la rigueur méthodologique exigée par son
statut de travail en cours. Les chapitres suivants
(Ch.~\ref{chap:cryptanal} à~\ref{chap:cas-zknr}) approfondissent
les outils de conception, de cryptanalyse et de vérification
formelle nécessaires pour, le jour où une publication complète
sur ZK-NR ou tout protocole similaire sera disponible, en conduire
vous-même l'audit rigoureux.
\end{boxinfo}